\section{Objetivos}

\subsection{Objetivo general}

Completar el caso de estudio 2 implementando concurrencia mediante \textbf{OpenMP} sobre la aplicación de multiplicación de matrices, caracterizar el comportamiento de la aplicación a nivel de CPU y de memoria, y medir su rendimiento en dos máquinas diferentes para poder comparar hardware.

\subsection{Objetivos específicos}

\begin{itemize}
    \item Implementar la versión concurrente con \textbf{OpenMP} partiendo del programa secuencial, aplicando las directivas de la API (\texttt{\#pragma omp}) de forma justificada.
    \item Justificar técnicamente la \textbf{elección de cada directiva y cláusula} utilizada, apoyándose en los ejemplos de la especificación oficial de OpenMP.
    \item Realizar el \textbf{perfilado mínimo de la aplicación} en dos ejes: \textbf{CPU} (tiempo por función, con GProf) y \textbf{memoria} (tasa de fallos de caché y footprint de memoria, con Cachegrind y \texttt{getrusage}).
    \item \textbf{Medir el rendimiento en dos máquinas distintas} con Linux, usando un benchmark reproducible, y comparar los resultados mediante la métrica de diferencia porcentual de rendimiento.
    \item Analizar de forma integrada los resultados de \textit{benchmarks}, \textit{profiling} y \textit{speedup} para explicar qué limita el rendimiento de la aplicación.
\end{itemize}
